% sec-02: Document 02, AB 3162, participant rights and robotic safety.
\docpart{AB 3162 Patient Rights}{AB 3162 - California LLM Oncology Patient
Rights and Robotic Safety Act of 2026}

\section{Findings and Declarations}

The Legislature finds and declares all of the following. The findings in
document 01 §1 of this package are incorporated and are not repeated.

\begin{enumerate}
\item A participant in an interventional oncology trial consents to an
  investigational agent and to an investigational procedure. Where a large
  language model also contributes to the conduct of that trial, the participant
  is entitled to know it, in the consent discussion and not only in the consent
  form \cite{cfr50adapt}.

\item Patient-facing work published in 2026 establishes that participants
  understand and use a stated right of refusal when it is offered in plain
  terms, and that the offer costs a site scheduling capacity rather than safety
  \cite{patientadvocacy,patientinstructions2026}.

\item A ten-gate assurance suite developed for a mobile surgical humanoid
  demonstrates that a robot in a clinical envelope can be held to a testable
  safety specification rather than to a manufacturer's assurance
  \cite{h2humanoid}.

\item Clinician trust in a model-assisted workflow is earned by supervision that
  is recorded, not by model accuracy alone, and the record of supervision is
  what a participant is entitled to rely on \cite{cliniciantrust}.

\item Existing California law addresses generative artificial intelligence in
  patient communications and in utilization review, but not the interventional
  trial setting \cite{ab3030_2024,sb1120_2024}.

\item An emergency stop that a participant's care team cannot reach within one
  step is not an emergency stop. Safety in a robotic procedure suite is a matter
  of measured distances and measured times, and this act states both.

\item A right that has no evidentiary record cannot be enforced. Each right
  created by this act is paired in §4 with the record that proves it was
  honored.

\item Remedies under this act are limited to injunctive relief, because a
  damages remedy in an investigational setting would deter the very sites the
  state is trying to authorize while adding nothing that the department's
  enforcement power under document 01 §5 does not already provide.
\end{enumerate}

\section{Definitions}

Participant, advisory-only output, robotic procedure suite, robot instance, and
site safety officer carry the meanings given in document 01 §2. In addition:

\begin{enumerate}
\item ``Human-only pathway'' means the conduct of a participant's trial
  procedures with no large language model output generated, retrieved, or
  displayed in connection with that participant's care.

\item ``Robotic operating envelope'' means the three-dimensional volume within
  which a robot instance may move during a study procedure, defined in document
  09 §3 and marked on the floor of the suite.

\item ``Emergency stop'' means a hardware-initiated command that arrests all
  robot motion, independent of software state.

\item ``Reportable robotic event'' means any unintended motion, any envelope
  breach, any interlock failure, and any emergency stop actuation, whether or
  not a participant was present.
\end{enumerate}

\section{Participant Rights}

A participant at a qualified site has all of the following rights, which may not
be waived by contract, by consent form, or by protocol.

\begin{enumerate}
\item The right to a human-only pathway, on request, at any time, without giving
  a reason and without effect on eligibility, on the assigned dose level, or on
  any other aspect of trial participation.

\item The right to be told, in the consent discussion, that a large language
  model contributes advisory output to the conduct of the trial, what it is used
  for, and what it is not permitted to do.

\item The right to know that no model output reaches the participant except
  through an affirmative act by a named clinician.

\item The right to a copy of the model output generated in connection with the
  participant's own care, and of the clinician's disposition of it, on request
  and at no charge.

\item The right to be informed of any reportable robotic event that occurred
  during a procedure in which the participant was present, within 72 hours.

\item The right to withdraw from the trial without prejudice to continuing care,
  and to a written statement of the care that will continue.

\item The right to raise a concern to the department directly, without going
  through the site, and to be told how.
\end{enumerate}

Subdivision (a) is stated first and without qualification because it is the
right that costs the site money. Its operational consequence is carried into
the standard operating procedure index at document 12 §3 as a named procedure
with a named owner.

\section{Disclosure and Consent}

\begin{mvtable}
\begin{tabularx}{\textwidth}{>{\raggedright\arraybackslash}p{4.0cm} >{\raggedright\arraybackslash}p{3.2cm} >{\raggedright\arraybackslash}X}
\headrow \thead{Disclosure} & \thead{Where it is made} & \thead{Evidence that it was made}\\
\midrule
A model contributes advisory output & Consent form and consent discussion & Signed consent form, and the coordinator's consent discussion note \\
What the model is not permitted to do & Consent discussion & The consent discussion note, which names the four prohibitions \\
The human-only pathway is available & Consent discussion and the participant handout & The handout acknowledgment line, initialed \\
The audit trail is retained & Consent form & The consent form section on records \\
A reportable robotic event occurred & Within 72 hours, in writing & The event notification letter and the delivery record \\
Withdrawal does not end clinical care & Consent form and at withdrawal & The withdrawal note and the continuing care statement \\
\end{tabularx}
\end{mvtable}
\tabcapl{tab:ab3162disc}{Disclosure, its place, and its evidence. Every row
names a record a monitor could ask for, because a disclosure obligation with no
record is unverifiable at a survey.}

The consent architecture follows the adapted human subject protection framework
rather than being invented here \cite{cfr50adapt,cfr50}. Where the adapted
framework and 45 CFR Part 46 differ, the more protective governs
\cite{cfr45part46}.

\section{Robotic Safety Requirements}

The numbers in this section are identical to document 09 §3, which is the
section that owns them. A site may not comply with one and not the other.

\begin{mvtable}
\begin{tabularx}{\textwidth}{>{\raggedright\arraybackslash}p{4.6cm} >{\raggedright\arraybackslash}p{3.0cm} >{\raggedright\arraybackslash}X}
\headrow \thead{Requirement} & \thead{Value} & \thead{Standard}\\
\midrule
Envelope boundary beyond maximum reach & 1.2 meters, floor-marked & Site specification \\
Emergency stop reach from any point inside the envelope & 0.8 meters & Site specification \\
Emergency stop at each suite door & Required & Site specification \\
Independent interlock channels & 3 & ISO 13849-1 category 3 \cite{iec80601277} \\
Motion arrest after a stop command & 250 milliseconds & IEC 80601-2-77 \cite{iec80601277} \\
Quasi-static force limit in collaborative operation & 80 newtons & ISO 13482 \cite{iso134822014} \\
Restart after an emergency stop & Two-person authorization & Document 09 §3 \\
\end{tabularx}
\end{mvtable}
\tabcapl{tab:ab3162safe}{The seven safety values. Each is a measurement a
surveyor can take with a tape, a stopwatch, or a force gauge, which is the test
this act applies to every safety claim.}

\section{Adverse Event Reporting and Remedies}

\begin{enumerate}
\item A reportable robotic event shall be recorded in the site's event log
  within one hour and reported to the department within 5 business days. An
  event resulting in participant injury shall be reported within 24 hours.

\item Adverse event reporting to the federal government is governed by the
  investigational new drug safety reporting requirements and is not displaced by
  this section \cite{cfr312,cfr312adapt}. Where the state and federal clocks
  differ, the shorter governs.

\item A participant, or the department on a participant's behalf, may bring an
  action for injunctive relief to enforce a right created by §3. This act
  creates no right to damages. The limit is stated in the same subdivision as
  the remedy so that it cannot be read as an oversight.
\end{enumerate}
